\section{Faktorisierung}

\begin{df}\label{skript:2.1}
	Sei $ f \in \H^2 $.
	Wir bezeichnen $ f $ als \textit{innere Funktion}, falls $ | f | = 1 $ fast überall auf $ \mathbb{T} $ gilt.
	Wir nennen $ f $ eine \textit{äußere Funktion}, falls
	\begin{align*}
	E_f := \bigvee_{n \in \N_0} \mathscr{S}^n f = \H^2
	\end{align*}
	gilt.
\end{df}

\begin{genericthm}{Faktorisierungstheorem}\label{skript:2.2}
	Sei $ f \in \H^2 \setminus \lbrace 0 \rbrace $.
	Dann lässt sich $ f $ bis auf eine multiplikative Konstante eindeutig als Produkt einer inneren Funktion $ f^i $ und einer äußeren Funktion $ f^a $ beschreiben.
	D.h. es gilt $ f = f^i f^a $ und $ E_f = f^i \H^2 $.
\end{genericthm}

\begin{proof}
	Wegen $ f \neq 0 $ erhalten wir  mit \ref{skript:1.3}
	\begin{align*}
	E_f = \Theta \H^2	
	\end{align*}
	mit $ |\Theta| = 1  $ f. ü. auf $ \mathbb{T} $.
	Damit ist $ \Theta  $ eine innere Funktion. Wir setzen $ f^i := \Theta $ und $ f^a := f \overline{\Theta} $, womit gilt $ f = f^i f^a $. 
	Wir müssen noch nachweisen, dass $ f^a $ eine äußere Funktion ist.
	Da $ f^i $ als Multiplikationsoperator eine Isometrie auf $ \L^2 $ bzw. $ \H^2 $ ist, erhalten wir 
	\begin{align*}
	E_f
	= \bigvee_{n \in \N_0} z^n f^i f^a 
	= f^i \bigvee_{n \in \N_0} z^n f^a 
	= f^i E_{f^a}.
	\end{align*}
	Damit folgt $ f^i H^2 = f^i E_{f^e} $ und damit auch $ H^2 = E_{f^e} $. Also ist $ f^e $ eine äußere Funktion.
	Die Eindeutigkeit erhalten wir aus Beuerling-Helson und \ref{skript:1.2}.
\end{proof}

\begin{genericthm}{Folgerung}
	Sei $ f \in \H^2$.
	Dann gilt:
	\begin{align*}
	f \ \text{ist äußere Funktion} \
	\Leftrightarrow
	\left(g \in \H^2, \frac{g}{f} \in \L^2 \Rightarrow \frac{f}{g} \in \H^2 \right)
	\end{align*}
\end{genericthm}

\begin{proof}
	Inhalt...
\end{proof}